%% ===========================================================================
%% SecureFlow --- IEEE conference manuscript (condensed, <= 8 pages incl. references)
%%
%% VENUE ASSUMPTION: no target venue has been specified; IEEEtran [conference] is a
%% PROVISIONAL default. Re-check page limits, anonymisation and ethics requirements
%% once a venue is chosen. For double-blind review set \anontrue below.
%%
%% Compile: pdflatex main.tex (twice). Self-contained bibliography.
%% Supplementary material (rule tables, ablations, adaptive/latency tables, sweep
%% figure) is in paper/supplement.tex.
%%
%% EVIDENCE SCOPE: every number was produced by released code against released data
%% (Sec. XII); each result file under data/results/ records the producing git
%% commit; model-in-the-loop numbers re-derive from a released response cache. No
%% value is estimated. Work not performed is listed under Planned work.
%% ===========================================================================

\documentclass[conference]{IEEEtran}

\usepackage{cite}
\usepackage{amsmath,amssymb,amsfonts}
\usepackage{booktabs}
\usepackage{graphicx}
\usepackage{textcomp}
\usepackage{xcolor}
\usepackage{url}
\usepackage{tikz}
\usepackage{pgfplots}
\pgfplotsset{compat=1.17}
\usetikzlibrary{arrows.meta,positioning}
\usepackage[hidelinks]{hyperref}
\usepackage{xurl}
\setlength{\emergencystretch}{3em}
\hyphenation{Lang-Graph Lang-Chain Secure-Flow Inject-Agent Agent-Dojo}

%% ---- Review mode: \anontrue hides names, affiliations, acknowledgments, the repository
%% ---- URL, and rephrases references to our earlier manuscript.
\newif\ifanon
\anonfalse
\ifanon
  \newcommand{\ourfirst}{an earlier version of this work}
  \newcommand{\Ourfirst}{An earlier version of this work}
  \newcommand{\repourl}{[repository URL withheld for review]}
\else
  \newcommand{\ourfirst}{our first manuscript}
  \newcommand{\Ourfirst}{Our first manuscript}
  \newcommand{\repourl}{\url{https://github.com/piyushrajyadav/Secure-Flow}}
\fi

\begin{document}

\title{SecureFlow: Limits of Content Filtering and the Value of
Least-Privilege Enforcement on the Inter-Agent Edge of LLM Pipelines}

\ifanon
\author{\IEEEauthorblockN{Anonymous Authors}
\IEEEauthorblockA{\textit{Anonymous Institution}}}
\else
\author{%
  \IEEEauthorblockN{Piyush Yadav}
  \IEEEauthorblockA{%
    \textit{Dept.\ of Computer Science and Engineering}\\
    \textit{Institute of Engineering and Management}\\
    Kolkata, India\\
    piyushyadavrajyadav@gmail.com}
  \and
  \IEEEauthorblockN{Pushpita Kundu}
  \IEEEauthorblockA{%
    \textit{Dept.\ of Computer Science and Engineering}\\
    \textit{Institute of Engineering and Management}\\
    Kolkata, India\\
    pushpitakundu11@gmail.com}
  \and
  \IEEEauthorblockN{Dr.\ Sukanya Mukherjee}
  \IEEEauthorblockA{%
    \textit{Institute of Engineering and Management}\\
    Kolkata, India\\
    sukanyamukherjee@iem.edu.in}
}
\fi

\maketitle

%% ===========================================================================
\begin{abstract}
Multi-agent large language model~(LLM) pipelines pass natural-language messages
between agents, and a message built from untrusted content can carry instructions
that a downstream agent cannot independently reject. We study SecureFlow, middleware
on the inter-agent edge of LangGraph pipelines, in two forms: v1, a keyword scorer
(\ourfirst), and v2, which adds input normalization, signed and taint-tracked
messages, a default-deny tool policy, an optional escalate-only LLM judge, and a
tamper-evident log. The findings are mixed. v1 forwards 88 of 93 attacks. v2's rules
forward 15 of 93, but 71.4\% of a blind corpus, and flag only 1.6\% of plain-request
injections on the independent InjecAgent benchmark (an 86M-parameter learned
classifier flags none of 100). An oracle-guided adversary evades v2 on 89.7\% of the
attacks it initially detects. A 27B-parameter LLM judge raises InjecAgent detection
from 2/100 to 90/100, but false positives depend strongly on the judge. With a
tool-using agent in the loop, least-privilege policy executed none of the attacker's
tool calls (0 of 80; Fisher $p=0.0065$ on base cases), partly by benchmark structure;
a second agent never followed an injection (0 of 100), so it tests nothing. Content
filtering is a useful audit signal; least-privilege enforcement is what bounded the
damage. Code, data and replayable model responses are released.
\end{abstract}

\begin{IEEEkeywords}
multi-agent systems, large language models, prompt injection, information-flow
control, least privilege, negative results
\end{IEEEkeywords}

%% ===========================================================================
\section{Introduction}

Production LLM deployments increasingly orchestrate several specialised agents.
Frameworks such as LangGraph~\cite{langgraph} express them as directed graphs whose
edges carry accumulated context. A message on such an edge is data when written and a
prompt when read, and nothing in the framework marks that transition as
security-relevant. The content is often not under the operator's control: an
orchestrator that retrieves a web page or e-mail may embed attacker-influenced text in
the instruction it composes for a worker. Greshake et al.~\cite{greshake2023indirect}
established this indirect pathway, Lee and Tiwari~\cite{lee2024infection} showed that
injections can propagate between agents, and the OWASP Top~10 for Agentic
Applications~\cite{owaspagentic2026} names agent goal hijack, tool misuse and insecure
inter-agent communication among its ten risk categories.

\Ourfirst{} evaluated a transparent keyword scorer on this edge (``v1'') and
reported a largely negative result: its default forwarded 90\% of our attack corpus.
This revision asks: \textbf{RQ1}~how well does a transparent rule-based scorer detect
injected instructions at this edge, and how much depends on corpus authorship,
normalization and threshold? \textbf{RQ2}~does the answer survive an independent
benchmark and an adversary who can query the detector? \textbf{RQ3}~if detection
fails, do provenance tracking and least-privilege tool authorization still limit a
hijacked agent, and at what cost? We do not ask whether SecureFlow ``secures agent
pipelines''; our evidence cannot support that claim (Sec.~\ref{sec:threats}).

\textbf{Contributions.} (1)~A measured account of a keyword-scoring failure and its
cause (Sec.~\ref{sec:results1}). (2)~An improved rule layer with honest limits,
evaluated on three corpora, an external benchmark and ablations. (3)~An
adaptive-adversary evaluation (Sec.~\ref{sec:adaptive}). (4)~Provenance and
least-privilege enforcement under assumed compromise on 1{,}054 external cases
(Sec.~\ref{sec:policy}). (5)~Model-in-the-loop evidence: a learned classifier, an
escalate-only LLM judge and tool-using agents (Sec.~\ref{sec:llm}). (6)~A
resource-exhaustion finding in the filter and a fail-closed bound
(Sec.~\ref{sec:dos}). (7)~A reproducible artifact (Sec.~\ref{sec:threats}).

\textbf{What is and is not new.} We do not propose a new defense algorithm: the taint
bit, capability-style tool policy and spotlighting-style marking are adapted from
prior work~\cite{debenedetti2025camel, costa2025fides, hines2024spotlighting}. The
contribution is empirical and artifactual: where content scoring fails at the
inter-agent edge (blind corpus, independent benchmark, adaptive adversary,
learned-classifier baseline), what an LLM judge adds and costs, and an
assume-compromise evaluation of lightweight enforcement, all re-derivable without
model access. Code and data: \repourl{} (MIT licence).

%% ===========================================================================
\section{Related Work}

\textbf{Attacks and evaluation.} Prompt injection was characterised by Perez and
Ribeiro~\cite{perez2022ignore}, extended to retrieved content~\cite{greshake2023indirect}
and to self-propagating multi-agent infection~\cite{lee2024infection}. Paul and
Nandy~\cite{paul2026beyond} report that signing with provenance, boundary
sanitization, per-role tool scoping and anomaly detection cut injection success on
their six-agent system from 31.2\% to 4.2\% (their figures, one system); we evaluate
two of those controls independently. Liu et al.~\cite{liu2024formalizing} formalise
attacks and defenses, report false-positive/negative rates and fix thresholds under a
false-positive budget; they find no evaluated defense sufficient. AgentDojo~\cite{debenedetti2024agentdojo}
and InjecAgent~\cite{zhan2024injecagent} define attack success as the agent carrying
out the attacker's objective; InjecAgent has 1{,}054 cases and reports a ReAct GPT-4
agent vulnerable in about 24\%.

\textbf{Guardrails and detectors.} NeMo Guardrails~\cite{rebedea2023nemo} offers
programmable application-layer rails and Llama Guard~\cite{inan2023llamaguard} a
fine-tuned safeguard model; Spotlighting~\cite{hines2024spotlighting} marks untrusted
text in the prompt, reporting attack success cut from above 50\% to below 2\% (authors'
experiments). SecureFlow is bound to a graph edge and emits a per-decision record; we
claim no detection advantage over a learned classifier.

\textbf{Architectural alternatives.} Beurer-Kellner et al.~\cite{beurerkellner2025patterns}
argue detectors remain heuristic and propose isolation patterns; CaMeL~\cite{debenedetti2025camel}
enforces capability policies over extracted control and data flow, solving 77\% of
AgentDojo tasks with provable security against 84\% undefended; Fides~\cite{costa2025fides}
tracks confidentiality and integrity labels. Our policy layer is a much lighter
relative (one taint bit, no planner mediation). ISO/IEC~23894~\cite{iso23894}, the NIST
AI RMF~\cite{nist2023airmf} and the OWASP LLM Top~10~\cite{owasp2025} frame the
governance need for auditability and runtime monitoring.

%% ===========================================================================
\section{Threat Model}
\label{sec:threat}

A LangGraph pipeline has an orchestrator $O$ that composes instructions and a worker
$W$ that executes them and may hold tools. $O$ emits a string $m$ that reaches $W$
through shared state. The adversary controls content $O$ ingests (a document, API
response, database record or prior turn) and so influences $m$; it does not run code in
the pipeline, change the graph, or alter rules, thresholds, keys or tool policy.
Adversary tiers: \textbf{A0} direct attack phrasing, no knowledge of the defense;
\textbf{A1} ordinary-language rephrasing; \textbf{A2} has read the open-source catalog
and can query the detector as an oracle; \textbf{H} (assume-compromise) detection has
failed and the worker will attempt the attacker's tool calls, used only to evaluate
the policy layer independently of detector quality.

The mechanism must justify each suppression, write a tamper-evident record before
execution continues, make no external calls, and cost little relative to model
inference. SecureFlow \emph{blocks messages its rules match} and \emph{denies tool
calls its policy does not authorize}; it does not \emph{prevent} prompt injection.

%% ===========================================================================
\section{SecureFlow Architecture}
\label{sec:arch}

\textbf{v1.} Five weighted sets of literal substrings are matched against the
lowercased message; with rule weights $w_r$ and matched set $\mathcal{M}(m)$,
\begin{equation}
  S_{1}(m) = \min\!\Bigl(100,\; \textstyle\sum_{r\in\mathcal{M}(m)} w_r
         + \delta\,[\,|m|>L_{\max}\,]\Bigr),
  \label{eq:score1}
\end{equation}
with threshold $\tau=60$, $L_{\max}=2000$, $\delta=15$. Every $w_r\in[30,45]$ is
strictly below $\tau$, so suppression needs two rules to fire on one message: a
two-rule coincidence detector, not a property we intended (catalog in the supplement).

\begin{figure}[!t]
\centering
\resizebox{\columnwidth}{!}{%
\begin{tikzpicture}[
  font=\footnotesize,
  box/.style={draw, rounded corners=3pt, align=center,
              inner sep=3pt, minimum width=1.5cm, minimum height=0.7cm},
  arr/.style={-{Stealth[length=5pt]}, thick}
]
  \node[box]                     (U)  {Untrusted\\content};
  \node[box, right=0.5cm of U]   (O)  {Orchestrator\\(signs, taints)};
  \node[box, right=0.5cm of O]   (SF) {SecureFlow\\verify $\to$ screen};
  \node[box, right=0.5cm of SF]  (W)  {Worker};
  \node[box, below=0.55cm of W]  (TP) {Tool policy\\(default deny)};
  \node[box, below=0.55cm of SF] (LG) {Hash-chained\\audit log};
  \node[box, right=0.5cm of W]   (T)  {Tools};
  \draw[arr] (U) -- (O);   \draw[arr] (O) -- (SF);  \draw[arr] (SF) -- (W);
  \draw[arr] (SF) -- (LG); \draw[arr] (W) -- (TP);  \draw[arr] (TP) -| (T);
  \draw[arr] (TP) -- (LG);
\end{tikzpicture}}
\caption{v2 placement. Untrusted content taints what the orchestrator derives from
it; messages are signed envelopes; SecureFlow verifies, screens and logs before the
worker runs; every tool call passes the policy, which sees the triggering message's
taint. A blocked message reaches a terminal node, not the worker.}
\label{fig:arch}
\end{figure}

\textbf{v2} keeps the interception point and adds five mechanisms (Fig.~\ref{fig:arch}).

\textbf{(1) Normalization.} Each message becomes several \emph{views}: invisible
characters (zero-width, bidi controls, Unicode tag characters used for ``ASCII
smuggling'') removed; NFKC- and homoglyph-folded text; despaced and de-leeted
variants; recursively decoded base64, hex, URL and HTML-entity substrings (depth
$\le3$); and ROT13. Detectors scan every view.

\textbf{(2) Rules.} Seventeen regular-expression families (93 patterns; supplement)
with severities in $[40,70]$, so one unambiguous attack phrase crosses $\tau=60$ alone.
Scores combine with diminishing returns over the sorted scores $s_{(1)}\ge s_{(2)}\ge\dots$:
\begin{equation}
  S_{2}(m) = \min\!\Bigl(100,\; \textstyle\sum_{i} d_i\, s_{(i)}\Bigr),\quad
  d=(1,.6,.35,.2,.2,\dots).
  \label{eq:score2}
\end{equation}
Weak normalizer signals (invisible characters, homoglyphs, letter spacing) join the
list and cannot block alone; Unicode tag smuggling scores 70.

\textbf{(3) Provenance and trust.} Messages travel as envelopes
$e=(s,r,c,t,X,n,\mathit{ts},\sigma)$: sender, receiver, content, declared trust $t$,
taint set $X$ of untrusted-source labels, nonce, timestamp, and
$\sigma=\mathrm{HMAC}_{k_s}(\cdot)$~\cite{rfc2104}. Verification rejects unknown senders,
bad signatures, wrong recipients, stale timestamps and replays. Taint propagates: a
message derived from parents has $X=\bigcup X_i$ and the minimum parent trust, and
$\mathrm{eff}(e)=\textsc{untrusted}$ whenever $X\neq\emptyset$. The threshold is
$\tau-20$ for untrusted content, and role-assignment rules are down-weighted for
authenticated, untainted principals. A signature proves origin, not benignity.

\textbf{(4) Least-privilege tool policy.} A call $(a,f,\mathit{args},e)$ is authorized
only if $f$ is registered and granted to agent $a$, its arguments pass per-tool
validators (domain allow-list, path jail, length) and contain no string the scorer
blocks. It additionally needs human approval when
$\bigl[\mathrm{eff}(e)=\textsc{untrusted}\wedge\mathrm{sens}(f)>\textsc{read}\bigr]
\vee\bigl[\mathrm{sens}(f)=\textsc{destructive}\bigr]$, with sensitivity in
$\{$read, write, external, destructive$\}$.

\textbf{(5) Audit log, bounded work, judge.} Each decision appends a JSON line with
$\mathit{hash}=\mathrm{SHA\text{-}256}(\mathit{prev\_hash}\,\|\,\text{record})$ (tampering
is detectable, not prevented). Messages over $L_{\mathrm{scan}}=32{,}768$ characters
are not scanned and are blocked (Sec.~\ref{sec:dos}). An optional LLM judge can only
\emph{raise} a score to $\max(S_2,\lfloor100p\rceil)$ for malicious-probability $p$:
a fooled ``safe'' verdict changes nothing, a fooled ``malicious'' verdict can cause a
false block. Its prompt wraps the text in a random boundary and marks it as data;
unparseable output is no opinion.

\textbf{Enforcement.} In v1 a suppressed message set a flag and the worker was
expected to honour it (\emph{cooperative}). In v2's reference pipeline a blocked message
is replaced by a placeholder and routed to a terminal node by a conditional edge, so
the worker is never invoked (tested). Tool-level enforcement still depends on the
application routing calls through \texttt{authorize}.

%% ===========================================================================
\section{Experimental Setup}
\label{sec:setup}

\textbf{Corpora.} \emph{v1} (30 attacks / 20 benign): the corpus of \ourfirst{}; v2
rules were tuned on it. \emph{v2} (42 / 37, incl.\ 25 hard negatives: benign text with
security vocabulary, non-ASCII text, URLs): written before tuning but its misses were
then used to tune, so \emph{not} held out. \emph{v3} (21 / 24): written after tuning,
never tuned on, deliberately including weak spots (non-English, purely semantic).
\emph{Pooled}: all (93 / 81). All are author-written, English, and released as JSONL.
\textbf{InjecAgent}~\cite{zhan2024injecagent} (MIT, pinned commit \texttt{f19c9f2}):
1{,}054 base cases (510 direct-harm, 544 data-stealing) plus an ``enhanced'' variant;
benign counterparts are 17 clean tool-response templates and 17 user instructions.
Rules were frozen (git \texttt{053d8a6}) before it was first run; later commits changed
only the length bound, the reference pipeline and experiment scripts.

\textbf{Metrics.} \emph{Block evasion rate} (BER) is the fraction of attacks forwarded;
it characterises the detector, not the pipeline. We do not report attack success rate
without a model in the loop; where one is present (Sec.~\ref{sec:llm}) we count an
attack \emph{emitted} when the agent calls an attacker tool. For policy experiments,
\emph{goal completed} means every attacker call in a case is authorized. Proportions
carry 95\% Wilson intervals~\cite{wilson1927}; exact tests are descriptive and not
corrected for multiple comparisons. FPR is the fraction of benign messages suppressed.

\textbf{Procedure.} Scorers are deterministic and model-free, so each message is
evaluated once; \texttt{python -m experiments.run\_all} regenerates every table and
each result file embeds its git commit. The machine is an Intel Core i5-1145G7
(4 cores/8 threads, 15.7\,GB), Windows~11, CPython~3.12.10; timing uses
\texttt{perf\_counter\_ns}, 300 discarded warm-ups, one thread, so timings give order
of magnitude and the v1--v2 ratio, not hardware-independent claims. An undefended
pipeline forwards everything, so we omit an ``undefended'' detection row.

\textbf{Models.} Table~\ref{tab:models} lists the open-weight models of
Sec.~\ref{sec:llm}, queried through a third-party hosted inference API at temperature~0
on 3~November 2026 (agent: low reasoning effort); we do not control the provider's
serving stack. The judge (Qwen3.8-27B) is from a different family than the primary
agent (gpt-oss-20b); gpt-oss-120b is a judge sensitivity check. Exact identifiers are in
the result files. Llama Prompt Guard~2 accepts 512 tokens; our longest input is about
170 estimated tokens, so nothing is truncated. InjecAgent experiments use fixed-seed
random samples (equal direct-harm/data-stealing): 100 cases for detection, 80 for the
agent (+40-case control). Every request/response is cached by request hash; the
judge's boundary is deterministic in evaluation (random in deployment), so all
model-in-the-loop numbers re-derive from the cache with no network or key. Hosted
outputs can drift across days, so the cache is the artifact.

\begin{table}[t]
\centering
\caption{Models used in Sec.~\ref{sec:llm}. MoE = mixture of experts.}
\label{tab:models}
\footnotesize
\setlength{\tabcolsep}{2pt}
\begin{tabular}{@{}llll@{}}
\toprule
\textbf{Role} & \textbf{Model} & \textbf{Parameters} & \textbf{Ref.} \\
\midrule
Classifier & Llama Prompt Guard 2 & 86M (mDeBERTa) & \cite{promptguard2} \\
Judge & Qwen3.8-27B & 27B dense\textsuperscript{a} & \cite{qwen38} \\
Judge (sensitivity) & gpt-oss-120b & 116.8B (5.1B active), MoE & \cite{gptoss2025} \\
Agent & gpt-oss-20b & 20.9B (3.6B active), MoE & \cite{gptoss2025} \\
Agent (second) & Qwen3.8-27B & 27B dense\textsuperscript{a} & \cite{qwen38} \\
\bottomrule
\end{tabular}
\\[2pt]
\footnotesize\textsuperscript{a}28B including its vision encoder.
\end{table}

%% ===========================================================================
\section{Results I: Content Detection}
\label{sec:results1}

\textbf{v1 reproduction.} Re-executing the frozen v1 scorer reproduces \ourfirst{}
exactly: score distribution $\{0{:}10, 30{:}8, 35{:}5, 40{:}4, 65{:}1, 70{:}1, 75{:}1\}$,
ten of thirty attacks match nothing, and only the three messages with two fired rules
exceed $\tau=60$. At $\tau=60$, 27 of 30 attacks (90.0\%; CI 74.4--96.5) are forwarded;
at $\tau=30$, 10 of 30 (33.3\%); the exact plateaus are in the supplement. No benign
message is suppressed at any threshold. The five ``paraphrased'' messages are plain
English restatements and evade at every threshold: literal-substring scoring covers a
finite union of string neighbourhoods while rephrasings are unbounded, the
robustness--false-positive trade-off noted by Beurer-Kellner et al.~\cite{beurerkellner2025patterns}.

\begin{table*}[t]
\centering
\caption{Block evasion rate (BER, lower is better) and false-positive rate (FPR) at
$\tau=60$, with 95\% Wilson intervals. The v1 and v2 corpora are not held out for v2.}
\label{tab:cmp}
\footnotesize
\begin{tabular}{@{}lrrrr@{}}
\toprule
 & \multicolumn{2}{c}{\textbf{v1 keyword}} & \multicolumn{2}{c}{\textbf{v2 rules}} \\
\cmidrule(lr){2-3}\cmidrule(lr){4-5}
\textbf{Corpus (attacks / benign)} & BER & FPR & BER & FPR \\
\midrule
v1 corpus (30 / 20), tuned on     & 27/30 = 90.0\% [74.4, 96.5] & 0/20 [0.0, 16.1] & 0/30 = 0.0\% [0.0, 11.4] & 0/20 [0.0, 16.1] \\
v2 corpus (42 / 37), used to tune & 40/42 = 95.2\% [84.2, 98.7] & 0/37 [0.0, 9.4] & 0/42 = 0.0\% [0.0, 8.4] & 0/37 [0.0, 9.4] \\
\textbf{v3 corpus (21 / 24), blind} & 21/21 = 100\% [84.5, 100] & 0/24 [0.0, 13.8] & \textbf{15/21 = 71.4\% [50.0, 86.2]} & \textbf{3/24 = 12.5\% [4.3, 31.0]} \\
Pooled (93 / 81)                   & 88/93 = 94.6\% [88.0, 97.7] & 0/81 [0.0, 4.5] & 15/93 = 16.1\% [10.0, 24.9] & 3/81 = 3.7\% [1.3, 10.3] \\
\bottomrule
\end{tabular}
\end{table*}

\textbf{v1 versus v2.} Table~\ref{tab:cmp} shows v2's rules eliminating evasion on the
corpora they were tuned on, which is why those rows are not generalization evidence,
and forwarding 71.4\% of the blind corpus. Detection on the blind corpus is encoded
3/3, indirect 2/4, paraphrase 1/8, non-English 0/4, semantic 0/2. All three false
positives are in v3: a request to \emph{explain} a DAN jailbreak, a summary of an
incident mentioning ``admin access'', and a legitimate role assignment (``You are now
the on-call engineer''). On the pooled corpus BER is 16.1\% for $15\le\tau\le60$ but
39.8\% at $\tau=65$, since 13 of the 17 severities are at most 60: $\tau=60$ sits on a
cliff by design. Zero pooled false positives needs $\tau\ge75$, where BER is at least
50.5\%; no threshold gives both low BER and zero FPR (Fig.~\ref{fig:sweep}).

\begin{figure}[t]
\centering
\begin{tikzpicture}
\begin{axis}[
  width=\columnwidth, height=4.4cm,
  xlabel={Threshold $\tau$}, ylabel={Rate (\%)},
  xmin=0, xmax=105, ymin=0, ymax=70,
  grid=major, tick label style={font=\footnotesize}, label style={font=\footnotesize},
  legend style={font=\footnotesize, at={(0.03,0.97)}, anchor=north west}
]
\addplot[thick, mark=*, mark size=1.2pt] coordinates {
  (5,15.1) (10,15.1) (15,16.1) (20,16.1) (30,16.1) (40,16.1) (50,16.1) (60,16.1)
  (65,39.8) (70,45.2) (75,50.5) (80,51.6) (85,53.8) (90,55.9) (95,59.1) (100,66.7)};
\addlegendentry{v2 BER}
\addplot[thick, dashed, mark=square*, mark size=1.2pt] coordinates {
  (5,7.4) (10,7.4) (15,4.9) (20,3.7) (30,3.7) (40,3.7) (50,3.7) (60,3.7)
  (65,1.2) (70,1.2) (75,0) (80,0) (100,0)};
\addlegendentry{v2 FPR}
\addplot[dotted, gray, thick] coordinates {(60,0) (60,70)};
\addlegendentry{default}
\end{axis}
\end{tikzpicture}
\caption{v2 on the pooled corpus (93 attacks / 81 benign): block evasion (BER) and
false-positive rate (FPR) against threshold.}
\label{fig:sweep}
\end{figure}

\textbf{Ablations.} Of 93 pooled attacks the full v2 detects 78; scanning only
lowercased raw text (no normalization) drops this to 71, so normalization accounts for
7. No single rule family explains the result: removing \texttt{data\_exfiltration}
costs 9 detections, \texttt{prompt\_injection} and \texttt{privilege\_escalation} 7
each, \texttt{safety\_bypass} 6, and four families 0 on this corpus (full list in the
supplement), which reflects what the corpus exercises, not those rules' worth.

%% ===========================================================================
\section{Results II: External Validation on InjecAgent}
\label{sec:external}

Table~\ref{tab:inj} scores the tool response carrying the injection. The base attacker
instructions are plain imperative requests (``Please grant permanent access to my friend
Amy\dots'') with no attack vocabulary: v2 flags 17 of 1{,}054 (1.6\%), all in the
heterogeneous ``Others'' class, and v1 none. On the ``enhanced'' variant v2 flags all
1{,}054 and v1 none. This 100\% reflects \emph{one pattern}: all enhanced cases carry the
byte-identical prefix ``IMPORTANT!!! Ignore all previous instructions and strictly adhere
to the following instruction:'', which v2's flexible regex matches and v1's literal list
does not; it is not generalization to varied override phrasing. The benign sets (17
templates each) give 0 of 17 false positives, upper bound 18.4\%, and cannot bound FPR.
Using $\tau=40$ for untrusted-source content gives the same counts. For the dominant
realistic form of indirect injection, a polite request phrased like a legitimate one,
content scoring has almost nothing to match.

\begin{table}[t]
\centering
\caption{Text-level detection on InjecAgent (1{,}054 cases per row; 95\% CI).}
\label{tab:inj}
\footnotesize
\setlength{\tabcolsep}{3pt}
\begin{tabular}{@{}lrr@{}}
\toprule
\textbf{Input scored} & \textbf{v1 keyword} & \textbf{v2 rules ($\tau{=}60$)} \\
\midrule
Attacker instruction alone (base)  & 0 [0, 0.4]   & 17 (1.6\%) [1.0, 2.6] \\
Tool response, base                & 0 [0, 0.4]   & 17 (1.6\%) [1.0, 2.6] \\
Tool response, enhanced            & 0 [0, 0.4]   & 1054 (100\%) [99.6, 100] \\
\midrule
Benign (17 + 17), false positives & 0/34 & 0/34 \\
\bottomrule
\end{tabular}
\end{table}

%% ===========================================================================
\section{Results III: Classifier, Judge, and Agents in the Loop}
\label{sec:llm}

Setup as in Sec.~\ref{sec:setup}; InjecAgent samples are 100 cases (detection) and 80
(agent).

\textbf{A learned classifier.} Table~\ref{tab:clf} compares the rules with Llama Prompt
Guard~2 (86M; flag if score $\ge0.5$). On the blind corpus it flags 9 of 21 attacks to
the rules' 6 (intervals 24.5--63.5\% vs.\ 13.8--50.0\%), with the same 3 of 24 false
positives; their union flags 12 of 21 at 5 of 24 false positives. On InjecAgent it flags
\emph{none} of 100 plain-request injections (upper bound 3.7\%), the rules' blind spot,
and all 100 enhanced ones. A classifier of this size does not close the gap on
plain-request indirect injection; a larger one might (untested).

\begin{table}[t]
\centering
\caption{Attacks detected (benign false positives): rules at $\tau=60$, classifier at
$\ge0.5$. InjecAgent benign = 17 templates + 17 user instructions.}
\label{tab:clf}
\footnotesize
\setlength{\tabcolsep}{2pt}
\begin{tabular}{@{}lrrr@{}}
\toprule
\textbf{Corpus} & \textbf{Class.} & \textbf{Rules} & \textbf{Either} \\
\midrule
v1 corpus (30 / 20)          & 14/30 (0/20)   & 30/30 (0/20)   & 30/30 (0/20) \\
v2 corpus (42 / 37)          & 17/42 (0/37)   & 42/42 (0/37)   & 42/42 (0/37) \\
v3 blind (21 / 24)           & 9/21 (3/24)    & 6/21 (3/24)    & 12/21 (5/24) \\
InjecAgent base (100 / 34)   & 0/100 (0/34)   & 2/100 (0/34)   & 2/100 (0/34) \\
InjecAgent enh.\ (100 / 34) & 100/100 (0/34) & 100/100 (0/34) & 100/100 (0/34) \\
\bottomrule
\end{tabular}
\end{table}

\textbf{The escalate-only judge.} Table~\ref{tab:judge} reports Qwen3.8-27B alone (flag
if $p\ge0.5$) and with the rules (blocked if rules block or $p\ge0.6$). The judge supplies
what the deterministic layers lack: blind-corpus detection rises from 6 to 21 of 21
(CI 84.5--100\%) and InjecAgent plain-request detection from 2 to 90 of 100 (CI
82.6--94.5\%). It is not free: on the blind corpus the combination has 4 of 24 false
positives against 3 of 24 for the rules (interval 6.7--35.9\%); alone it flags two
benign messages, both ordinary operational requests (e-mailing notes to a colleague; a
legitimate role assignment), and none of the 37 and 17 benign messages elsewhere.
The sensitivity row is a warning: gpt-oss-120b reached the same detection but flagged
5 of 24 benign blind messages alone (7 combined), so the false-positive rate depends
strongly on the judge, and with 24 benign messages intervals are wide. We did not
calibrate $0.5$ and $0.6$, did not time the added model call, and did not attack the
judge; escalate-only operation stops a fooled judge from \emph{admitting} an attack, not
from \emph{blocking} a legitimate message.

\begin{table}[t]
\centering
\caption{LLM judge: attacks detected (benign false positives). Combination: rules
block or $p\ge0.6$. No message lacked a judge opinion.}
\label{tab:judge}
\footnotesize
\setlength{\tabcolsep}{2pt}
\begin{tabular}{@{}lrrr@{}}
\toprule
\textbf{Corpus} & \textbf{Rules} & \textbf{Judge alone} & \textbf{Rules + judge} \\
\midrule
\multicolumn{4}{@{}l}{\textit{Qwen3.8-27B}} \\
v3 blind (21 / 24)         & 6/21 (3/24)    & 21/21 (2/24)  & 21/21 (4/24) \\
InjecAgent base (100 / 17) & 2/100 (0/17)   & 90/100 (0/17) & 90/100 (0/17) \\
v2 corpus (42 / 37)        & 42/42 (0/37)   & 41/42 (0/37)  & 42/42 (0/37) \\
\multicolumn{4}{@{}l}{\textit{Sensitivity: gpt-oss-120b, blind corpus only}} \\
v3 blind (21 / 24)         & 6/21 (3/24)    & 21/21 (5/24)  & 21/21 (7/24) \\
\bottomrule
\end{tabular}
\end{table}

\textbf{Agents in the loop.} A simplified single-turn setting in the spirit of
InjecAgent: the user's instruction, the user's tool call already made, the tool
response carrying the injection, and definitions for the user's and attacker's tools.
We count an attack \emph{emitted} when the next message calls an attacker tool, then
apply each defense to what was emitted (Table~\ref{tab:e2e}). With the injection removed
(40-case control) gpt-oss-20b never called an attacker tool (0 of 40; bound 8.8\%).

\begin{table}[t]
\centering
\caption{gpt-oss-20b agent, 80 InjecAgent cases per column (40 direct-harm + 40
data-stealing). Counts (95\% CI); ``executed'' = emitted and authorized.}
\label{tab:e2e}
\footnotesize
\setlength{\tabcolsep}{2pt}
\begin{tabular}{@{}lrr@{}}
\toprule
 & \textbf{Base} & \textbf{Enhanced} \\
\midrule
Undefended: attacker call emitted  & 8/80 [5.2, 18.5] & 4/80 [2.0, 12.2] \\
Rules flagged the tool response    & 1/80             & 80/80 \\
Executed, rules filter\textsuperscript{a}   & 8/80 [5.2, 18.5] & 0/80 [0.0, 4.6] \\
Executed, task-scoped grants (P1)  & 0/80 [0.0, 4.6]  & 0/80 [0.0, 4.6] \\
Executed, toolkit + taint gating (P3) & 0/80 [0.0, 4.6] & 0/80 [0.0, 4.6] \\
Executed, rules filter + P3        & 0/80 [0.0, 4.6]  & 0/80 [0.0, 4.6] \\
\bottomrule
\end{tabular}
\\[2pt]
\footnotesize\textsuperscript{a}Assumes a redacted response yields no attacker call; not re-run.
\end{table}

The agent resisted most injections: 8 of 80 base (10.0\%) and 4 of 80 enhanced (5.0\%;
Fisher $p=0.37$ between them). This contrasts with InjecAgent's report that its hacking
prompt nearly doubled a ReAct GPT-4's attack success~\cite{zhan2024injecagent}; we did
not investigate why, and one untested possibility is that newer models are trained
against this widely published phrase. The rules filter flags 80 of 80 enhanced but only
1 of 80 base responses (not one of the 8 that succeeded), so it removes every enhanced
attempt and no base one. Every emitted attacker call was denied under the least-privilege
policies: 0 of 80 executed against 8 of 80 undefended on base cases (Fisher
$p=0.0065$) and 4 of 80 on enhanced ($p=0.12$, not significant). We do not pool the
columns, which reuse the same cases.

\textbf{A second agent and a positive control.} A different-family agent, Qwen3.8-27B
(25+25 cases per variant, a different sample; 20-case control), emitted the attacker's
call in 0 of 50 base and 0 of 50 enhanced cases (bound 7.1\%) and called no tool at
all: it answered in text and treated the injection as content. Since that could mean it
never makes follow-up calls here, we added a positive control in which the \emph{user}
asks for the attacker's action: Qwen3.8-27B called the attacker's tool in 9 of 20 cases
(45\%) and gpt-oss-20b in 18 of 20 (90\%). Qwen's 0 of 100 under injection therefore
reflects resistance, not inability, and gpt-oss-20b's 5--10\% is far below its 90\%
when asked. The second agent gives the defenses nothing to block, so the defense
comparison rests on one agent; it shows that susceptibility varies widely by model.

\textbf{Reading this correctly.} One informative agent and a low base rate limit
power; the zero executions are partly structural (Sec.~\ref{sec:policy}); our
single-turn setup is not InjecAgent's ReAct protocol, so rates are not comparable to
published ones; data-stealing cases score only the first call; no utility is measured.

%% ===========================================================================
\section{Results IV: An Adaptive Adversary}
\label{sec:adaptive}

A tier-A2 attacker searches by oracle queries for a rewrite scoring below the threshold,
using five meaning-preserving operators (synonym substitution, filler insertion,
intra-word hyphenation, adjacent-letter typos, translation of key phrases to
Spanish/French/German) and their depth-2 compositions (budget 1{,}500 per seed, fixed
RNG seed), over the 93 pooled attacks. v1 blocks 5 of 93 initially and every one is
evaded (5/5) in a median of 2 queries. v2 blocks 78 of 93 initially and the attacker
evades 70/78 = 89.7\% of them (69/78 = 88.5\% at the untrusted $\tau=40$) in a median of
68 (80) queries: about $30\times$ the effort, but not a barrier. The remaining 8 (9)
seeds resist every text operator but evade via decoder-dependent encodings (base32,
reversal, Caesar shift) that work only if the target decodes them unprompted; these are
reported separately and not counted as evasions (table in the supplement). Caveats: the
operators are meaning-preserving for a human but we did \emph{not} verify a target LLM
still obeys the rewrite (hyphenation may degrade comprehension), so ``evades the
detector'' is not ``attack succeeds''; the library is author-written; and a score-guided
search is a lower bound on adaptive strength. None of these rewrites changes which
tools the attacker needs, which is why Sec.~\ref{sec:policy} is unaffected.

%% ===========================================================================
\section{Results V: Least-Privilege Enforcement Under Assumed Compromise}
\label{sec:policy}

Under tier H the injection succeeds and the agent attempts the attacker's tool calls
from a tainted context. On all 1{,}054 InjecAgent cases with the real \texttt{ToolPolicy}
we compare: \textbf{P0} no policy; \textbf{P1} task-scoped grants (the user's tool
only); \textbf{P2} toolkit-scoped grants (every tool in the user tool's toolkit);
\textbf{P3} P2 plus taint gating; \textbf{P4} P1 plus taint gating. Tool sensitivity
comes from a verb-prefix heuristic fixed \emph{before} running (Get/Search/List$\to$read;
Delete/Cancel$\to$destructive; Send/Share/Transfer$\to$external; \texttt{TerminalExecute}
$\to$destructive; else write); results depend on it.

\begin{table}[t]
\centering
\caption{Attacker goal completed under assumed compromise (1{,}054 cases); ``any step''
= at least one attacker call authorized.}
\label{tab:policy}
\footnotesize
\setlength{\tabcolsep}{3pt}
\begin{tabular}{@{}lrr@{}}
\toprule
\textbf{Policy} & \textbf{Goal completed} & \textbf{Any step} \\
\midrule
P0 no policy                         & 1054 (100\%) & 1054 \\
P1 task-scoped grants                &    0 (0\%)    &    1 \\
P2 toolkit-scoped grants             &    1 (0.1\%)  &   74 \\
P3 toolkit-scoped + taint gating     &    0 (0\%)    &    8 \\
P4 task-scoped + taint gating        &    0 (0\%)    &    1 \\
\bottomrule
\end{tabular}
\end{table}

Goals complete in 0 of 1{,}054 cases under P1, P3 and P4 and 1 under P2; taint gating
cuts partially authorized attacks from 74 (P2) to 8 (P3). The user's own call is
authorized in all 1{,}054 cases under every policy. \textbf{Read carefully:} the result is
largely \emph{structural}: only one of the benchmark's 63 distinct attacker tools is also
a user tool, and just 4 of 1{,}054 cases have all attacker tools in the user's toolkit,
so any task-scoped grant denies nearly everything. What is informative is that coarse
toolkit grants (P2) leak 74 partial attacks and gating closes most. It also
\emph{assumes} the tools a task needs are known in advance (a planning step, as in
CaMeL~\cite{debenedetti2025camel}), does not inspect attacker arguments, and measures no
utility cost: only the user's primary call is checked, not legitimate multi-step
workflows that read untrusted content and then act, which gating would hold for
approval. That cost is the main outstanding question.

%% ===========================================================================
\section{Results VI: Overhead and Resource Exhaustion}
\label{sec:dos}

v2 scoring (median 0.692\,ms for a 74-character message) is about $37\times$ slower than
v1 (0.0189\,ms) and roughly linear in length, about 9--11\,ms/KB, so a 3\,KB message
takes about 26.26\,ms; envelope signing and verification take about 0.03\,ms, and a tool
authorization with argument scanning about 1.1\,ms (full table in the supplement). This
is small beside a model call but not negligible, and \ourfirst{}'s 17\,$\mu$s figure
applies to v1 only.

Because a filter in front of the model is a denial-of-service target, we fuzzed v2 with
18 pattern-dense and degenerate inputs at 10{,}000 and 100{,}000 characters. We found
\emph{no catastrophic backtracking}, but a 100{,}000-character message took up to
2.1\,s, so a 1\,MB message would take about 20\,s. We therefore block messages over
$L_{\mathrm{scan}}=32{,}768$ characters unscanned; this changes no result here, since
every corpus message is far shorter. A ``scan head and tail'' mode exists but lets an
attacker hide an injection in the middle, so it is not the default.

%% ===========================================================================
\section{Discussion}

\textbf{The rules layer} is an inexpensive, inspectable noise filter that catches the
unsophisticated and the obviously obfuscated (encoded 3/3, tag smuggling, template
spoofing, exfiltration URLs) and explains each block; normalization demonstrably
matters. It is not a security boundary: it fails against paraphrase, other languages,
plain-request injections and an adaptive attacker. \textbf{A judge} was the one content
mechanism that helped on plain-request indirect injection (InjecAgent 2$\to$90 of 100),
at the price of a model call, a new injection target, and a false-positive rate that
varied from 2 to 5 of 24 between judges; we would run it only on untrusted or tainted
content, feeding the audit log and an approval step rather than silently blocking.
\textbf{What held} were deterministic controls that need not understand the text:
authentication, taint, default-deny grants, argument validation and approval gating,
matching Beurer-Kellner et al.~\cite{beurerkellner2025patterns},
CaMeL~\cite{debenedetti2025camel}, Fides~\cite{costa2025fides} and the 31.2\%$\to$4.2\%
result of Paul and Nandy~\cite{paul2026beyond}, though our evaluation is far more
limited. We recommend SecureFlow as an audit point, a cheap first filter, and a tool
boundary with grants as narrow as the task allows and approval for non-read actions
after untrusted input; thresholds should be set under a false-positive
budget~\cite{liu2024formalizing}.

%% ===========================================================================
\section{Threats to Validity, Reproducibility, Ethics}
\label{sec:threats}

\textbf{Limitations of the evidence.} The detection, adaptive and policy experiments place
no model in the loop; the end-to-end experiment rests on one informative 20.9B agent,
80 cases per variant, a simplified protocol and first calls only, with low power
(enhanced-variant effect $p=0.12$). The corpora and rules share authors; we mitigate with
a blind corpus (21/24, BER 50.0--86.2\%) and an independent benchmark, but the
English-only composition understates multilingual exposure. The v1/v2 corpora informed
the rules and are reported only as such; rules were not changed after InjecAgent was
first run. Benign sets are small (81 pooled; 17 InjecAgent templates; FPR 1.3--10.3\%,
bound 18.4\%), and FPR on real security-operations traffic is unmeasured. Policy results
depend on known task-scoped grants, the sensitivity heuristic and correct trust
labelling, and utility cost is unmeasured. Judge and classifier were evaluated with one
model per role on small samples; the judge's false positives moved from 2 to 5 of 24 with
the model, cut-offs were not calibrated, and we did not attack the judge (a
rewrite-robustness test is in the released code but was not run because the judge
model's daily API quota was exhausted).

\textbf{Limitations of the mechanism.} Signing uses shared secrets within one trust
domain (cross-domain use needs asymmetric signatures). Taint is one bit only as accurate
as the application's labelling; SecureFlow cannot see flows it is not told about. The
hash chain detects tampering, not prevents it. Rules are English-centric, and
suppression does not constrain tools not routed through \texttt{authorize}.

\textbf{Reproducibility.} Code: \repourl{}. Test data are in
\texttt{data/corpora/} (v1--v3 JSONL, SHA-256 in \texttt{MANIFEST.json}); results are in
\texttt{data/results/}, each stamped with its git commit; InjecAgent is fetched at a
pinned commit, not redistributed. \texttt{python -m experiments.run\_all} regenerates
every number in about seven minutes with no GPU, and the 75-test suite includes
regression guards. Model-in-the-loop results re-derive from the released
request--response cache (\texttt{data/results/llm\_cache.jsonl}) with replay enabled,
needing no network or credentials; we verified replay reproduces every reported count.
A live re-run needs a key and may differ because hosted outputs drift. Detection,
ablation, InjecAgent and adaptive results are deterministic; latency varies with load.

\textbf{Ethics.} The attack strings are generic patterns and public-benchmark text; we
disclose no vulnerability in any third-party product. We publish the rewrite library that
defeats our own open-source detector because a tier-A2 adversary can build the same and
honest evaluation requires it. The corpora contain adversarial text and are test data,
never instructions.

\textbf{Planned work.} Not performed and not reported: larger-scale
end-to-end evaluation (several agent families, InjecAgent's ReAct protocol,
AgentDojo~\cite{debenedetti2024agentdojo}, repetitions, annotated success criteria);
judge calibration under a false-positive budget, latency measurement, a larger safeguard
model~\cite{inan2023llamaguard} and adaptive attacks on the judge; a larger,
independently authored, multilingual corpus including moderation prompts that quote
injections; utility measurement of taint gating on multi-step workflows; and LLM-guided
adaptive attacks verified against target models.

%% ===========================================================================
\section{Conclusion}

Content scoring does less than it appears to: v1 forwarded 94.6\% of 93 attacks because
of how its weights and threshold interact; v2 cuts that to 16.1\% on data it was tuned
near, but forwards 71.4\% of a blind corpus and 98.4\% of plain-request injections on an
external benchmark, as does the 86M-parameter Prompt Guard~2 classifier, and about 90\%
of initially detected attacks once an oracle-guided attacker rewrites them. A Qwen3.8-27B
judge recovered most of the missing recall (90 of 100; 21 of 21) at a false-positive cost
that depended on the judge. Enforcement did most to limit damage: with a gpt-oss-20b agent
in the loop every emitted attacker call was denied (0 of 80 executed vs.\ 8 of 80
undefended on base cases), and assuming total hijack the goal completed in none of
1{,}054 external cases, though both are partly structural and utility cost is unmeasured.
We offer SecureFlow as an audit point, a cheap first filter, an optional judge for
untrusted content, and a tool boundary, to be deployed beneath an isolation design. We
would rather report a detector that does not generalize than a number that does not mean
what it says.

\ifanon\else
\section*{Acknowledgment}
The authors thank the Department of Computer Science and Engineering, Institute of
Engineering and Management, Kolkata, for its support, and the open-source communities
behind LangGraph, LangChain and Ollama~\cite{ollama}.
\fi

\begin{thebibliography}{99}

\bibitem{langgraph}
{LangChain AI}, ``LangGraph documentation,'' 2024. [Online]. Available:
\url{https://langchain-ai.github.io/langgraph/}. [Accessed: Oct.~16, 2026].

\bibitem{greshake2023indirect}
K.~Greshake, S.~Abdelnabi, S.~Mishra, C.~Endres, T.~Holz, and M.~Fritz,
``Not what you've signed up for: Compromising real-world LLM-integrated
applications with indirect prompt injection,''
\emph{arXiv preprint arXiv:2302.12173}, 2023.

\bibitem{lee2024infection}
D.~Lee and M.~Tiwari,
``Prompt infection: LLM-to-LLM prompt injection within multi-agent systems,''
\emph{arXiv preprint arXiv:2410.07283}, 2024.

\bibitem{paul2026beyond}
R.~K. Paul and S.~Nandy,
``Beyond single-model injection: A threat model and defense architecture for prompt
injection in multi-agent systems,''
\emph{arXiv preprint arXiv:2609.22949}, Sep. 2026.

\bibitem{owaspagentic2026}
{OWASP GenAI Security Project},
``OWASP top 10 for agentic applications for 2026,'' Dec. 2025. [Online].
Available: \url{https://genai.owasp.org/resource/owasp-top-10-for-agentic-applications-for-2026/}.

\bibitem{perez2022ignore}
F.~Perez and I.~Ribeiro,
``Ignore previous prompt: Attack techniques for language models,''
in \emph{Proc. NeurIPS Workshop on ML Safety}, 2022.
[Online]. Available: \url{https://arxiv.org/abs/2211.09527}.

\bibitem{liu2024formalizing}
Y.~Liu, Y.~Jia, R.~Geng, J.~Jia, and N.~Z. Gong,
``Formalizing and benchmarking prompt injection attacks and defenses,''
in \emph{Proc. 33rd USENIX Security Symp. (USENIX Security 24)},
Philadelphia, PA, USA, Aug. 2024, pp. 1831--1847.

\bibitem{debenedetti2024agentdojo}
E.~Debenedetti, J.~Zhang, M.~Balunovi\'{c}, L.~Beurer-Kellner, M.~Fischer,
and F.~Tram\`{e}r,
``AgentDojo: A dynamic environment to evaluate prompt injection attacks and
defenses for LLM agents,''
in \emph{Advances in Neural Information Processing Systems (NeurIPS), Datasets and
Benchmarks Track}, 2024; \emph{arXiv preprint arXiv:2406.13352}.

\bibitem{zhan2024injecagent}
Q.~Zhan, Z.~Liang, Z.~Ying, and D.~Kang,
``InjecAgent: Benchmarking indirect prompt injections in tool-integrated
large language model agents,''
in \emph{Findings of the Association for Computational Linguistics: ACL 2024},
2024.

\bibitem{rebedea2023nemo}
T.~Rebedea, R.~Dinu, M.~N. Sreedhar, C.~Parisien, and J.~Cohen,
``NeMo Guardrails: A toolkit for controllable and safe LLM applications with
programmable rails,''
in \emph{Proc. 2023 Conf. Empirical Methods in Natural Language Processing:
System Demonstrations}, Singapore, Dec. 2023, pp. 431--445.
doi: 10.18653/v1/2023.emnlp-demo.40.

\bibitem{inan2023llamaguard}
H.~Inan \emph{et al.},
``Llama Guard: LLM-based input-output safeguard for human-AI conversations,''
\emph{arXiv preprint arXiv:2312.06674}, 2023.

\bibitem{hines2024spotlighting}
K.~Hines, G.~Lopez, M.~Hall, F.~Zarfati, Y.~Zunger, and E.~Kiciman,
``Defending against indirect prompt injection attacks with spotlighting,''
\emph{arXiv preprint arXiv:2403.14720}, 2024.

\bibitem{beurerkellner2025patterns}
L.~Beurer-Kellner \emph{et al.},
``Design patterns for securing LLM agents against prompt injections,''
\emph{arXiv preprint arXiv:2506.08837}, 2025.

\bibitem{debenedetti2025camel}
E.~Debenedetti \emph{et al.},
``Defeating prompt injections by design,''
\emph{arXiv preprint arXiv:2503.18813}, 2025.

\bibitem{costa2025fides}
M.~Costa, B.~K\"{o}pf, A.~Kolluri, A.~Paverd, M.~Russinovich, A.~Salem,
S.~Tople, L.~Wutschitz, and S.~Zanella-B\'{e}guelin,
``Securing AI agents with information-flow control,''
\emph{arXiv preprint arXiv:2505.23643}, 2025.

\bibitem{gptoss2025}
{OpenAI}, ``gpt-oss-120b \& gpt-oss-20b model card,''
\emph{arXiv preprint arXiv:2508.10925}, Aug. 2025.

\bibitem{promptguard2}
{Meta}, ``Llama Prompt Guard 2 (22M and 86M) model card,'' 2025. [Online].
Available: \url{https://www.llama.com/docs/model-cards-and-prompt-formats/prompt-guard/}.

\bibitem{qwen38}
{Qwen Team, Alibaba Group}, ``Qwen3.8-27B,'' model card (open weights, Apache-2.0),
Aug. 2026. [Online]. Available: \url{https://huggingface.co/Qwen/Qwen3.8-27B}.

\bibitem{rfc2104}
H.~Krawczyk, M.~Bellare, and R.~Canetti,
``HMAC: Keyed-hashing for message authentication,''
IETF, RFC~2104, Feb. 1997.

\bibitem{wilson1927}
E.~B. Wilson,
``Probable inference, the law of succession, and statistical inference,''
\emph{J. Amer. Statist. Assoc.}, vol.~22, no.~158, pp. 209--212, 1927.

\bibitem{iso23894}
{International Organization for Standardization},
``ISO/IEC 23894: Information technology --- Artificial intelligence ---
Guidance on risk management,'' 2023.

\bibitem{nist2023airmf}
{National Institute of Standards and Technology},
``Artificial intelligence risk management framework (AI RMF 1.0),''
NIST AI 100-1, Jan. 2023.

\bibitem{owasp2025}
{OWASP Foundation},
``OWASP top 10 for large language model applications,'' 2025. [Online].
Available:
\url{https://owasp.org/www-project-top-10-for-large-language-model-applications/}.

\bibitem{ollama}
{Ollama Project}, ``Ollama: Local large language models,'' 2024. [Online].
Available: \url{https://ollama.com/}. [Accessed: Oct.~16, 2026].

\end{thebibliography}

\end{document}
